% from ega0-19-fr.tex
\oldpage[0\textsubscript{IV-1}]{101}

donc comme $C$-module à gauche (M, IX, 4). Or, comme $K$ est une extension
finie séparable de $k$, on sait que $K\otimes_k K$ est composée directe
d'extensions de $k$, dont l'une est $K$ lui-même (Bourbaki, \emph{Alg.},
chap.~VIII, \textsection8, prop.~3)~; cela entraîne donc que $K$ est un
$C$-module \emph{projectif}, d'où notre assertion. Toute $k$-extension de $K$
de noyau $L$ est donc $k$-triviale, et \emph{a fortiori} les $k$-extensions
commutatives le sont, d'où le théorème dans ce cas (19.4.4).

\emph{II. --- Cas général.} --- Avec les notations de (18.4.5), il s'agit de
prouver que $\operatorname{Hom}_K(\mathrm H_2(P'_\bullet),L)=0$ pour
\emph{tout} $K$-espace vectoriel $L$, ce qui signifie évidemment que
$\mathrm H_2(P'_\bullet)=0$. Si $K$ est réunion d'une famille filtrante de
sous-extensions $K^{(\alpha)}$ de $k$, $P'_\bullet$ est \emph{limite
inductive} des complexes correspondants $P'^{(\alpha)}_\bullet$, puisque le
foncteur $\varinjlim$ commute au produit tensoriel dans la catégorie des
$k$-modules~; par l'exactitude du foncteur $\varinjlim$ dans cette catégorie,
on a donc
$\mathrm H_2(P'_\bullet)=\varinjlim\mathrm H_2(P'^{(\alpha)}_\bullet)$.
Comme l'hypothèse que $K$ est séparable entraîne qu'il en est de même de toute
sous-extension de $K$ et que $K$ est réunion des sous-extensions de type fini,
la première partie de la démonstration entraîne bien que $K$ est une
$k$-algèbre formellement lisse. C.Q.F.D.
\end{proof}

\begin{corollary}[19.6.2]
\label{0.19.6.2-fr}
Soient $A$ un anneau local séparé et complet, $K$ son corps résiduel, $k$ un
sous-corps de $K$ tel que $K$ soit une extension séparable de $k$. Alors il
existe un sous-corps $K'$ de $A$ contenant $k$, tel que la restriction à $K'$
de l'homomorphisme canonique $A\to K$ soit un isomorphisme de $K'$ sur $K$.
\end{corollary}

\begin{proof}
Soit $\mathfrak m$ l'idéal maximal de $A$. En vertu de l'hypothèse et de
(19.6.1), $K$ est une $k$-algèbre formellement lisse~; appliquant (19.3.10) en
remplaçant $C$ par $A$ et $\mathfrak J$ par $\mathfrak m$, on voit que
l'automorphisme identique de $K=A/\mathfrak m$ se factorise en
$K\xrightarrow{u}A\to A/\mathfrak m$, où $u$ est un $k$-homomorphisme, donc
nécessairement un $k$-isomorphisme de $K$ sur un sous-corps $K'$ contenant
$k$.
\end{proof}

\begin{corollary}[19.6.3]
\label{0.19.6.3-fr}
Soient $A$ un anneau local noethérien complet, $\mathfrak m$ son idéal maximal,
$k$ son corps résiduel. Les conditions suivantes sont équivalentes~:
\begin{enumerate}
  \item[a)] $A$ contient un sous-corps.
  \item[b)] Si $p$ est la caractéristique de $k$, on a $pA=0$.
  \item[c)] Il existe un corps $K$ tel que $A$ soit isomorphe à un anneau
  quotient d'un anneau de séries formelles $B=K[[T_1,\ldots,T_n]]$.
\end{enumerate}
Lorsqu'il en est ainsi, on peut supposer que $A$ est isomorphe à $B/\mathfrak b$,
où $\mathfrak b$ est contenu dans le carré de l'idéal maximal de $B$.
\end{corollary}

\begin{proof}
Il est immédiat que $c)$ implique $a)$, car $A$ est $\neq0$ puisque c'est un
anneau local~; comme $K$ est un corps et que l'homomorphisme composé
$K\xrightarrow{\varphi}B\to A$ (où $\varphi$ est l'injection canonique) n'est
pas nul, il est injectif. Pour voir que $a)$ entraîne $b)$, il suffit de
remarquer que si $K$ est un sous-corps de $A$, $K$ est isomorphe à un
sous-corps de $k$ et a par suite même caractéristique $p$~; comme $p.1=0$ dans
$K$, donc dans $A$, on a $pA=0$. Réciproquement, $b)$ entraîne $a)$, car le
composé des homomorphismes canoniques
$\mathbf Z\xrightarrow{j}A\to k$ a pour noyau $p\mathbf Z$, et l'on a
$j(p\mathbf Z)=0$, donc $\operatorname{Ker}(j)=p\mathbf Z$ et $A$ contient un
anneau $A'$ isomorphe à $\mathbf Z/p\mathbf Z$ et ne rencontrant pas
$\mathfrak m$~; si $p>0$, $A'$ est un corps~; sinon, tout

